\section{Exclusion of every $G_2$ and triality normal core}
\label{sec:g2-triality-cores}

\paragraph{Status and input boundary.}
This new author derivation has passed an independent whole-proof and
source audit by \texttt{/root/g34\_eightfolds}.  That reviewer also
reconstructed every original generic entry and every small-cover spectrum
with separate exact arithmetic implementations; both passed.  The
root reviewer has separately passed the complete mathematical and code
review and rerun every certificate; its remaining source read is
recorded separately.  The complete published
character tables, their official generic-table implementation, and the
published list of exceptional Schur multipliers are mathematical inputs.
The new bounds below are proved from those inputs by rational polynomial
inequalities and exact cyclic Fourier calculations.  In particular, the
assertions for unbounded field size do not extrapolate a finite scan.
The source log distinguishes original papers actually retrieved from
published inputs accessed through their official mathematical tables.

As before, for a finite-order operator on an $a$-dimensional complex space,
write
\[
 \nu(T)=a-\max_{\lambda}\dim\ker(T-\lambda).
\]
Multiplication by a scalar does not change $\nu$.  A genuine nonidentity
scalar inertia has ordinary moving rank $a$.

\subsection*{D-G2T-01. Uniform bounds from the complete generic tables}

\begin{lemma}[D-G2T-01: all-field character bounds]
\label{D-G2T-01}
Let $S$ be $G_2(q)$ with $q\ge5$, or ${}^3D_4(q)$ with $q\ge3$.
For every nontrivial ordinary irreducible character $\chi$ of $S$ and
every $s\ne1$,
\[
                  |\chi(s)|\le\tfrac14\chi(1).             \tag{G2T.1}
\]
Moreover $\chi(1)\ge124$ in the first family and
$\chi(1)\ge219$ in the second.
\end{lemma}

\begin{proof}
The complete ordinary $G_2$ tables are the published inputs of
Enomoto--Yamada in characteristic two, Enomoto in characteristic three,
and Chang--Ree in the other characteristics
\cite{G2T:EY,G2T:Enomoto,G2T:ChangRee}.
For triality we use Deriziotis--Michler's complete table
\cite[\S4, Theorem~4.3 and Tables~3.6, 4.4]{G2T:DM}.
The official GenericCharacterTables implementation, descended from
CHEVIE, records these exact sources and supplies all the mathematical
entries \cite{G2T:GCT,G2T:CHEVIE}.
We freeze its source at commit
\texttt{87231738271ed548f62e7f1c98ff2bbfc1ed30d3}.

The following inventory is exhaustive.  The integer in the fourth
column is the number of character types and also the number of
class types; a type may contain several parameterized characters
or classes.  The last column is a proved lower bound for every
nontrivial character degree throughout the stated real range.
\[
\begin{array}{c|l|r|r|r}
\text{table}&\text{field condition}&q_0&\text{types}&\text{degree bound}\\\hline
\texttt{G2.01}&q\text{ even},\ q\equiv1\ (3)&5&27&126\\
\texttt{G2.02}&q\text{ even},\ q\equiv2\ (3)&5&27&124\\
\texttt{G2.101}&q=3^f,\ q\equiv1\ (4)&5&28&280\\
\texttt{G2.103}&q=3^f,\ q\equiv3\ (4)&5&28&280\\
\texttt{G2.111}&q\text{ odd},\ q\equiv1\ (3),\ q\equiv1\ (4)&5&31&126\\
\texttt{G2.113}&q\text{ odd},\ q\equiv1\ (3),\ q\equiv3\ (4)&5&31&126\\
\texttt{G2.121}&q\text{ odd},\ q\equiv2\ (3),\ q\equiv1\ (4)&5&31&124\\
\texttt{G2.123}&q\text{ odd},\ q\equiv2\ (3),\ q\equiv3\ (4)&5&31&124\\
\texttt{3D4.0}&q\text{ even}&3&29&219\\
\texttt{3D4.11}&q\text{ odd},\ q\equiv1\ (4)&3&33&219\\
\texttt{3D4.13}&q\text{ odd},\ q\equiv3\ (4)&3&33&219
\end{array}                                                    \tag{G2T.2}
\]
For each type, all admissible values of every class and character
parameter are retained in the published interpretation.  Our estimates
are independent of their values, as follows.

Every table entry is a finite expression
\[
                z(q)=\sum_\alpha p_\alpha(q)u_\alpha,
                \qquad p_\alpha(q)\in\mathbf Q[q],
                \qquad |u_\alpha|=1.                       \tag{G2T.3}
\]
Here each $u_\alpha$ is a product of the roots of unity occurring
in that entry.  Equal formal products may be collected, while
relations between different products need not be used.  The
original generic cyclotomics have exponents that are rational
expressions in the real integral parameters, so each phase
indeed has absolute value one for every admissible specialization.

Put $x=q-q_0\ge0$ and expand exactly
\[
 p_\alpha(q_0+x)=\sum_k c_{\alpha k}x^k,
 \qquad d(q_0+x)=\sum_k d_kx^k,
\]
where $d$ is the degree polynomial of the row.  The triangle inequality
gives the uniform bound
\[
                 |z(q)|\le
                 \sum_k\left(\sum_\alpha|c_{\alpha k}|\right)x^k.
                                                               \tag{G2T.4}
\]
Thus it suffices to check, for every nontrivial row and every
nonidentity column,
\[
                  d_k-4\sum_\alpha|c_{\alpha k}|\ge0
                         \quad\hbox{for every }k.           \tag{G2T.5}
\]
Every one of these rational inequalities is nonnegative in all
eleven tables.  Also every $d_k$ is nonnegative, and $d_0$ is at
least the number recorded in (G2T.2).  This proves the claimed
uniform degree bounds as well as (G2T.1).

Here are the complete finite verification details.
\texttt{build\_generic\_input.py} reads the eleven frozen source
files, checks their separately recorded SHA-256 hashes, and
extracts only the original character entries, declared degrees,
and declared class lengths.  It does not execute downloaded
Julia code.  The resulting offline mathematical input
\texttt{generic\_character\_entries.json} has $415124$ bytes and
SHA-256
\texttt{66ad9b13d457d141d424084336787b3bb1f9d53818475b806f7a6cc24eb3664b}.

\texttt{parse\_generic\_bounds.py} accepts only integer constants,
the indeterminate $q$, rational constant division, addition,
multiplication, nonnegative integer powers, and formal unit phases.
It checks that the trivial row consists of ones, that every first
column entry equals the separately declared degree, and that the
first class length is one while all other class lengths are
strictly greater than one throughout the relevant range.  The
last assertions follow from nonnegative translated coefficients
and a positive constant term after subtracting one.
It then constructs and checks (G2T.5) for all $318$ nontrivial
row types and all $9242$ nontrivial-row, nonidentity-column
entries.  All checks pass.

The complete rational coefficient lists, excluding run-specific timing
data, are retained in
\texttt{generic\_polynomial\_bounds.json}; their SHA-256 is
\texttt{568425a2cf62e62dbe24709a3372d9bd182946e0a53ae74e7fed45ba640dcd13}.
A second, short program, \texttt{check\_generic\_polynomials.py},
uses only Python's exact rational arithmetic to recheck every
shifted coefficient inequality and every minimum degree from
these lists.  It also passes.  This second program validates
the polynomial certificate; the first program separately
verifies its correspondence with the retained original table
expressions.  Neither is offered as a new proof of the published
character-table classification itself.
\end{proof}

\subsection*{D-G2T-02. Exceptional covers and all outer normalizers}

\begin{lemma}[D-G2T-02: normalizer bounds]
\label{D-G2T-02}
Let $N$ be a finite quasisimple group acting faithfully and
irreducibly in degree $a$, whose simple quotient is $G_2(q)$,
$q\ge3$, or ${}^3D_4(q)$, $q\ge2$.
Every nonscalar finite-order element of its complex linear
normalizer satisfies
\[
                         \nu(T)\ge3a/16.                  \tag{G2T.6}
\]
Every such $a$ is at least twelve.  If the simple quotient is
$G_2(3)$, $G_2(4)$, or ${}^3D_4(2)$, the stronger bound
$\nu(T)\ge a/4$ holds and the respective minimum degrees are
$14,12,26$.  In the other cases the degree bounds in
D-G2T-01 apply.
\end{lemma}

\begin{proof}
The generic simply connected groups of these two types have
trivial center.  The only exceptional multipliers in these
families are those of $G_2(3)$ and $G_2(4)$; the complete
published exceptional list is \cite[Chapter~5, Table~1,
numbered p.~35]{BI:HHM}.  In particular there is no omitted
triality central cover.  For the three small simple groups,
the universal covers and their full cyclic outer extensions
are the following complete official CTblLib inputs:
\[
\begin{array}{c|r|r|r|r|c}
\text{table}&\text{classes}&\text{selected rows}&
 \text{nonscalar spectra}&\min a&\min\nu/a\\\hline
\texttt{3.G2(3)}&45&44&1848&14&3/7\\
\texttt{3.G2(3).2}&39&20&740&14&3/7\\
\texttt{2.G2(4)}&55&54&2862&12&1/3\\
\texttt{2.G2(4).2}&65&48&3024&12&1/3\\
\texttt{3D4(2)}&35&34&1156&26&4/13\\
\texttt{3D4(2).3}&49&39&1872&26&4/13
\end{array}                                                    \tag{G2T.7}
\]
The extension metadata record respectively multiplier and outer
orders $(3,2),(2,2),(1,3)$, in the sense of CTblLib
\S3.7-3 \cite{SPOR:CTblLib}.  Every source string records
the ATLAS; the group identities and complete tables are
published mathematical inputs.

The GAP exporter retains every character, class size, element
order, prime power map, source string, and actual core-to-extension
class fusion.  It verifies full character-row orthogonality.
The standalone verifier selects every nonlinear core character
and every extension character whose restriction equals a
nonlinear irreducible core character, by exact equality of
all its character values.  It checks fusion orders, trivial
identity kernel, normal-class size unions, and compatibility
with all retained prime power maps.

For an element $s$ of order $o$, every eigenvalue multiplicity is
computed by the exact Fourier formula
\[
       m_b=\frac1o\sum_{j=0}^{o-1}\chi(s^j)\zeta_o^{-bj},
                      \quad 0\le b<o.                    \tag{G2T.8}
\]
Prime power maps compute every power class.  Cyclotomic
polynomials and integer coefficient lists perform all arithmetic;
every $m_b$ is checked to be a nonnegative integer and their
sum is checked to equal $a$.  Scalar classes are recognized
by $|\chi(s)|^2=a^2$, which is equivalent to scalar action for
a finite-order unitary operator.  The full computation comprises
$239$ selected rows and $11502$ nonscalar spectra, with precisely
the minima in (G2T.7).  It passes, and every ratio is at least
one quarter.

The input \texttt{g2\_triality\_small\_inputs.json} has SHA-256
\texttt{e4a8aefbf5a4a32a5b31b217143ccae03db8b9f3753cc2e91a6169d0364b6a7c}.
The author programs are \texttt{export\_g2\_triality\_small.g}
and \texttt{check\_g2\_triality\_small.py}.  The latter reuses
the fully reviewed exact cyclotomic Fourier helper from D-SPOR.

We explain why these six fixed tables cover arbitrary compatible
small-group normalizers, including scalar changes in a power
relation.  Lift the representation to the universal perfect cover
$U$.  Every automorphism of its simple quotient has a unique
lift to $U$.  An inner normalizer is scalar times a represented
element of $U$.  For an outer normalizer, choose a matching
generating coset in the recorded extension $A=U.r$,
where $r=2$ or $3$.  After an inner adjustment, its action on
$U$ agrees with that of an element $t\in A$.  Write $t^r=u_0$.
Schur's lemma gives $T^r=c\rho(u_0)$ for a nonzero scalar $c$.
After multiplying $T$ by an $r$th root of $c^{-1}$, its
conjugation and power relations give a representation of this
fixed $A$ extending $\rho$.  It is irreducible because its
restriction to $U$ is irreducible.  Thus its row is among the
tested rows.  Nonfaithful central kernels cause no loss: all
nonlinear rows of the universal cover were retained.  This proves
the small-group quarter bound with no presumption that every
outer automorphism preserves every central kernel.

For all other field sizes, D-G2T-01 gives the character ratio
$1/4$ on the centerless core.  The complex commutator argument
of D-BI-18 gives (G2T.6) for every normalizer, not merely for
inner elements.  Explicitly, if $T$ is nonscalar, there is
$u\in N$ with $[T,u]\ne1$, by irreducibility.  The rank
inequality
\[
                \operatorname{rank}([T,u]-1)\le2\nu(T)
\]
holds after subtracting a largest-eigenspace scalar from $T$.
If $[T,u]$ is nonscalar, its unitary trace and (G2T.1) imply
\[
 a-2\operatorname{rank}([T,u]-1)
       \le\operatorname{Re}\chi([T,u])\le a/4,
\]
so $\nu(T)\ge3a/16$.  A nonidentity scalar commutator instead
has full moving rank $a$, giving the stronger $\nu(T)\ge a/2$.
The smallest degrees are those already certified in (G2T.2)
and (G2T.7), proving every assertion.
\end{proof}

\subsection*{D-G2T-03. The whole finite-cover source exclusion}

\begin{theorem}[D-G2T-03: $G_2$ and triality core exclusion]
\label{D-G2T-03}
Let an inducing block of a finite-cover Weil sixfold source have
finite image $J\subset\operatorname{GL}(V)$.  Suppose $J$ has
an irreducibly acting normal central product of quasisimple
groups whose simple quotients belong to
\[
                G_2(q)\ (q\ge3),\qquad
                {}^3D_4(q)\ (q\ge2).
\]
Then the actual period rank of the source is at most eight.

Mixed factor types, arbitrary scalar phases, compatible outer
automorphisms, and factor permutations are allowed.  No
primitivity, tensor-indecomposability, essential-self-duality,
equality of inducing and original fields, or nonsplit-polarization
hypothesis is required.
\end{theorem}

\begin{proof}
Write the normal central product as $N=N_1\cdots N_t$.
Its irreducible representation has the form
$V=V_1\otimes\cdots\otimes V_t$, with faithful irreducible
factor representations of dimensions $a_i\ge12$.
Set $e=\dim V=\prod_i a_i$.  The simple direct factors of
$N/Z(N)$ are intrinsic.  Each $N_i$ is the derived subgroup
of the inverse image of its simple factor, so $J$ permutes
the $N_i$.  Their generated full matrix algebras
$\operatorname{End}(V_i)\otimes1$ are consequently permuted
as well.  Every normalizing operator therefore consists of
factor intertwiners and a permutation of the appropriate
isomorphic factor spaces.

If the operator preserves all factors and is nonscalar,
one factor intertwiner is nonscalar.  Fixing all other
eigenvalues shows that the largest eigenspace fraction of
the tensor operator is no larger than the largest such
fraction in that factor.  Formula (G2T.6) therefore gives
$\nu(g)\ge3e/16$.

For a nontrivial permutation cycle of length $l\ge2$ on
factors of equal dimension $a$, tensor changes of basis
put its action in the form of cyclic word rotation with
a single diagonal twist.  The product twist is diagonalizable:
the finite order of the whole operator implies finite order
up to scalar for its powers on each tensor cycle.
Each orbit of the resulting basis words contributes at most
one vector to a specified eigenspace.  There are $a$ constant
words, and every other orbit has size at least two.  Hence
the largest eigenspace dimension is at most
\[
                    a+(a^l-a)/2\le3a^l/4.
\]
Tensoring with the other cycles preserves this bound, so
$\nu(g)\ge e/4\ge3e/16$.  This is the weighted cyclic-word
argument of D-SPOR-03 and does not assume that the twists
have distinct eigenvalues.  Nonidentity scalar inertia has
ordinary moving rank $e$.  Thus every genuine inertia has
ordinary moving rank at least $3e/16$.

Suppose the actual period rank were nine.  Apply the proved
intermediate-cover factorization D-BI-01 to this inducing
block.  It remains valid when the inducing character field
is larger than the original quadratic field: the first
complex Hodge filtration factors through that intermediate
cover, and the conjugate filtration is its Hermitian
annihilator.  If the intermediate base has genus $h$ and
$R$ genuine branch points, the multiplicity-six Euler formula is
\[
                      D=6+2e(1-h),\qquad R\le D.           \tag{G2T.9}
\]
For $h\ge2$, its right side is negative because $e\ge12$.
For $h=1$, the actual period map factors through a finite
Hurwitz cover of dimension at most six.  Consequently rank
nine requires $h=0$ and $R\ge12$, and then $D=2e+6$.
The inertia estimate proves the necessary inequality
\[
              12\left\lceil\frac{3e}{16}\right\rceil
                   \le2e+6,
                   \qquad e\le24.                        \tag{G2T.10}
\]
This is an obstruction obtained from factorization of the
actual variation, not a sufficiency inference from a source
dimension count.

If $t\ge2$, then $e\ge12^2=144$, contradicting (G2T.10).
For a single factor outside the three small cases,
D-G2T-01 gives $e\ge124$ or $e\ge219$, also impossible.
For a single factor in those cases, the quarter bound of
D-G2T-02 gives $3e\le2e+6$, hence $e\le6$, contrary to
the respective minima $14,12,26$.  Every case is excluded.
Since the target Weil period domain has dimension nine,
the actual period rank is at most eight.
\end{proof}

\paragraph{Boundary of the result.}
The nonsimple group $G_2(2)$ is not included.  Its derived
simple group is $U_3(3)$, as recorded in
\cite[Chapter~5, numbered p.~35]{BI:HHM}, and belongs to a
separate unitary-group gate.  No assertion about the other
classical or exceptional families is inferred from this
calculation.
